\documentclass[11pt,a4paper]{article}
\usepackage[top=0.65in, bottom=0.65in, left=0.8in, right=0.8in]{geometry}
\usepackage[T1]{fontenc}
\usepackage[utf8]{inputenc}
\usepackage{enumitem}
\usepackage{hyperref}
\usepackage{microtype}
\usepackage{titlesec}

\titleformat{\section}{\large\bfseries}{}{0em}{}[\rule{\textwidth}{0.5pt}]
\titlespacing{\section}{0pt}{8pt}{4pt}

\setlist{nosep, leftmargin=*}
\pagestyle{empty}
\hypersetup{colorlinks=false}

\begin{document}

\begin{center}
  {\LARGE \textbf{Dr. Rachel Green, MD}}\\[3pt]
  Internal Medicine \quad|\quad Board Certified\\[4pt]
  Boston, MA \quad|\quad
  \href{mailto:rachel.green.md@email.com}{rachel.green.md@email.com} \quad|\quad
  (617) 555-0284
\end{center}

\section{Medical Education}
\noindent\textbf{Doctor of Medicine (M.D.)}, Harvard Medical School, Boston, MA \hfill \textit{June 2015}\\
AOA (Alpha Omega Alpha Honor Society) \quad|\quad Research Distinction\\[3pt]
\noindent\textbf{Residency, Internal Medicine}, Massachusetts General Hospital, Boston, MA \hfill \textit{2015 -- 2018}\\
Chief Resident, PGY-3\\[3pt]
\noindent\textbf{Fellowship, Hospital Medicine}, Brigham and Women's Hospital, Boston, MA \hfill \textit{2018 -- 2019}\\[3pt]
\noindent\textbf{B.S. Biology (Pre-Med)}, Yale University, New Haven, CT \hfill \textit{May 2011}\\
Summa Cum Laude \quad|\quad GPA: 3.97/4.0

\section{Clinical Experience}

\noindent\textbf{Attending Physician -- Hospital Medicine} \hfill \textit{Aug 2019 -- Present}\\
\textit{Massachusetts General Hospital, Boston, MA}
\begin{itemize}
  \item Provide inpatient care for complex medical patients on 24-bed general medicine service; supervise 2 residents and 2 medical students per rotation.
  \item Manage high-acuity cases including sepsis, pulmonary embolism, acute coronary syndrome, and decompensated heart failure.
  \item Serve as attending of record for rapid response and code blue events; RRT response time consistently below hospital average.
  \item Quality improvement project: implemented standardized discharge checklist reducing 30-day readmission rate by 12\% (published in JHOSPMED 2022).
  \item Clinic: Maintain panel of 280 primary care patients; same-day sick visit access rate: 94\%.
\end{itemize}

\section{Licensure \& Certifications}
\textbf{Medical License:} Massachusetts (Active), California (Active)\\
\textbf{Board Certification:} American Board of Internal Medicine (ABIM) -- Internal Medicine, 2018 (Recertified 2024)\\
\textbf{DEA:} Active (Schedule II--V)\\
\textbf{BLS, ACLS, PALS:} Current\\
\textbf{USMLE:} Step 1: 265, Step 2 CK: 272, Step 3: 260

\section{Publications}
\begin{enumerate}[leftmargin=*, label={[\arabic*]}]
  \item \textbf{R. Green}, T. Singh, A. Kim. ``Standardized Discharge Protocols Reduce 30-Day Readmissions in General Medicine.'' \textit{J Hosp Med.} 2022;17(8):582--588.
  \item \textbf{R. Green}, et al. ``Troponin Trends and Outcomes in Non-ACS Hospitalized Patients.'' \textit{Am J Med.} 2021.
\end{enumerate}

\section{Teaching}
\begin{itemize}
  \item Core faculty, HMS Clerkship (3rd-year Internal Medicine); rated 4.9/5.0 by students (2022, 2023)
  \item Lecturer: MGH Grand Rounds on ``Diagnostic Uncertainty in Complex Inpatients'' (2023)
\end{itemize}

\section{Professional Memberships}
American College of Physicians (ACP, Fellow) \quad|\quad Society of Hospital Medicine (SHM) \quad|\quad AMA

\end{document}
